%!TEX root = ../hutbthesis_main.tex
% 设置中文摘要
% 说明：本摘要为开题阶段占位框架，实验数据、指标结论须在系统实现并测试通过后据实补写，未测试通过的功能与数据不得写入论文。
\keywordscn{水下仿真，传感器建模，水声通信链路，声呐成像，延迟与噪声，操作训练}
\begin{abstractzh}

海洋工程、水下探测、水下搜救与水下作业训练等场景中，水下环境具有噪声大、信号延迟高、载体漂移明显、视觉遮挡严重、水体能见度低等特点，光学传感器与水声传感器的感知质量受到显著影响，同时水声通信链路存在带宽窄、多径效应强、传播时延大、丢包率高等问题，真实环境下数据采集成本高、风险大、重复实验困难。本文围绕水下工作环境的传感器与通信链路仿真展开研究，基于仿真平台构建水下传感器与通信链路仿真系统，对水下噪声、延迟、漂移、遮挡、低能见度、声呐成像等环境因素进行建模与可视化呈现，并对传感器数据与通信链路的传输过程进行模拟，为感知算法、通信策略与水下操作训练提供可重复、低成本的实验环境。

本文主要研究内容包括：水下工作环境建模与环境参数配置；光学/水声等传感器在噪声、延迟、漂移、遮挡、低能见度条件下的感知仿真；水声通信链路的延迟、带宽、丢包与多径等特性仿真；声呐成像结果的模拟与显示；以及面向算法验证和操作训练的系统交互与测试评估。相关功能将在指定仿真场景下完成功能测试与性能测试，实验结果与量化指标在系统实现后据实补充。

\end{abstractzh}
